\section{The Missionary Journeys}

Acts organizes Paul's Aegean-and-Anatolian career into three great
circuits radiating from Antioch. The letters know nothing of this tidy
scheme---they reveal a man in nearly constant motion, with plans that
changed under pressure (2~Cor 1:15--17)---and Acts' three journeys should
be held as Luke's true-to-life but authorial arrangement, not as Paul's
own itinerary-book. Within that caution, the outline the two witnesses
share is firm: from Antioch through Cyprus and south-central Asia Minor;
then across to Macedonia and Achaia; then a long residence at Ephesus;
with Jerusalem at the hinge points.

\subsection{Antioch, the sending church}

Antioch on the Orontes, third city of the empire, is where the mixed
mission first became institutional: refugees of the Stephen persecution
preached there to Greeks as well as Jews, Barnabas was sent from Jerusalem
to inspect, fetched Saul from Tarsus, and for ``a whole year'' they taught
``a great multitude, so that at Antioch the disciples were first named
Christians'' (Acts 11:26). It was the Antiochene assembly, Acts says, that
heard the call---``the Holy Ghost said to them: Separate me Saul and
Barnabas, for the work whereunto I have taken them'' (13:2)---and released
the two with fasting and imposition of hands. Paul's letters confirm
Antioch's centrality from the other direction: it is the scene of his
confrontation with Cephas (Gal 2:11), which presupposes that it was his
own base and that its table-practice mattered to everyone.

\subsection{The first journey: Cyprus and Galatia (c.~46--48)}

The first circuit (Acts 13--14) runs Seleucia--Cyprus (Salamis to Paphos,
where the proconsul Sergius Paulus believes and the narrative quietly
switches from ``Saul'' to ``Paul,'' 13:9--12)---Perga---Antioch of
Pisidia---Iconium---Lystra---Derbe, and back through the same cities to
the Syrian Antioch. Two scenes carry the theology. At Pisidian Antioch,
after the synagogue sermon and its rejection: ``To you it behoved us
first to speak the word of God: but because you reject it\ldots\ behold we
turn to the Gentiles'' (13:46)---the pattern, synagogue first, then the
nations, that Acts will repeat in city after city and that Paul states as
principle: ``to the Jew first, and to the Greek'' (Rom 1:16). At Lystra,
after a healing, the crowds cry ``The gods are come down to us in the
likeness of men,'' calling Barnabas Jupiter and Paul Mercury ``because he
was chief speaker'' (14:10--11)---and then, swayed by rivals, stone Paul
and drag him out ``thinking him to be dead'' (14:18). The letters supply
the flat corroboration: ``once I was stoned'' (2~Cor 11:25); ``I bear the
marks of the Lord Jesus in my body'' (Gal 6:17). Whether the churches of
this circuit (Pisidian Antioch, Iconium, Lystra, Derbe, in the Roman
province Galatia) are the addressees of the letter to the Galatians---the
``South Galatian'' hypothesis---or whether that letter went to ethnic
Galatia further north, visited later (Acts 16:6; 18:23), is a standing
critical question that this study leaves open; it affects the letter's
date by several years but not the shape of the life.

\subsection{The pattern of the mission}

Before the second circuit, the pattern deserves statement, because it, and
not any single journey, is the historical achievement. Paul worked cities,
not countryside---provincial capitals and ports on the Roman roads and sea
lanes. He arrived at the synagogue, where there were Scriptures, Gentile
sympathizers, and a hearing; he stayed with, and worked alongside,
tradespeople like Aquila and Priscilla, ``working night and day, lest we
should be chargeable to any of you'' (1~Thess 2:9); assemblies met in
households; and when he moved on he left behind co-workers and wrote
letters. The mission was corporate throughout: Barnabas, Silvanus/Silas,
Timothy, Titus, Luke and the ``we''-narrator, Apollos; the deacon Phoebe,
``our sister, who is in the ministry of the church, that is in Cenchrae''
(Rom 16:1); Prisca ranked with her husband; Lydia the purple-dealer at
Philippi (Acts 16:14--15); Epaphras, Epaphroditus, Aristarchus, Mark,
``I Tertius, who wrote this epistle'' (Rom 16:22). Romans 16 alone greets
some two dozen named persons, a third of them women. Paul defended the
whole apparatus as apostolic right freely waived---``Have we not power to
carry about a woman, a sister, as well as the rest of the apostles, and
the brethren of the Lord, and Cephas?'' (1~Cor 9:5)---and stated its
governing flexibility as principle: ``I became to the Jews, a Jew, that I
might gain the Jews\ldots\ I became all things to all men, that I might
save all'' (1~Cor 9:20, 22). The signs and
sufferings were equally structural: ``Of the Jews five times did I receive
forty stripes, save one. Thrice was I beaten with rods, once I was stoned,
thrice I suffered shipwreck\ldots\ in perils of robbers\ldots\ in perils
from false brethren'' (2~Cor 11:24--26)---a catalogue, written before
Acts, that guarantees how much traveling and how much violence the tidy
circuits compress.

\subsection{The second journey: into Macedonia and Achaia (c.~49--52)}

After the council and a falling-out with Barnabas over Mark (Acts
15:37--40), Paul took Silas overland through Syria, Cilicia, and the
first-journey churches, adding Timothy at Lystra. Acts then narrates the
turn of the mission into Europe with a vision at Troas: ``a man of
Macedonia standing and beseeching him\ldots\ Pass over into Macedonia, and
help us'' (16:9)---and the narrative shifts into ``we.'' At Philippi, a
Roman colony, the first named European converts: Lydia and her household,
a slave-girl freed of a divining spirit, the jailer; also flogging and the
stocks, answered the next morning with citizenship: ``They have beaten us
publicly, uncondemned, men that are Romans'' (16:37). At Thessalonica,
weeks of synagogue argument, a riot, and a night evacuation; 1~Thessalonians,
written months later from Corinth, remembers the sequel from inside: ``you
turned to God from idols, to serve the living and true God'' (1~Thess 1:9).
Beroea received the word eagerly; Athens heard the Areopagus speech, Luke's
masterpiece of gospel-to-Greeks---``I found an altar also, on which was
written: To the unknown God. What therefore you worship, without knowing
it, that I preach to you'' (17:23)---and answered mostly with irony
(``some indeed mocked,'' 17:32), though ``Dionysius, the Areopagite, and a
woman named Damaris'' believed (17:34). Benedict~XVI's chronology
catechesis rightly treats the scene as programmatic: ``the discourse of
the Areopagus\ldots\ is the model of how to translate the Gospel into
Greek culture, of how to make Greeks understand that this God of the
Christians and Jews was not a God foreign to their culture but the
unknown God they were awaiting'' (27 August 2008). Then Corinth: eighteen months
(18:11), the trade with Aquila and Priscilla, the household baptisms, the
hearing before Gallio that anchors the chronology, and departure by
Cenchrae ``into Syria'' (18:18). The Thessalonian correspondence belongs
in this window; its burning question---the fate of believers who die
before the Lord's coming---received the answer that ``the dead who are in
Christ, shall rise first'' (1~Thess 4:15, DR numbering).

\subsection{The third journey: Ephesus and the crisis years (c.~52--57)}

The last circuit is really a residence: Ephesus, metropolis of Asia, for
``three months'' in the synagogue, then two years in the hall of Tyrannus,
``so that all they who dwelt in Asia, heard the word of the Lord'' (Acts
19:8--10); Paul himself: ``I will tarry at Ephesus until Pentecost. For a
great door and evident is opened unto me: and many adversaries'' (1~Cor
16:8--9). From Ephesus he governed his Aegean churches by courier and
letter: the Corinthian correspondence (at least four letters, of which two
survive whole), the fight with Corinth that included a ``painful visit''
and a letter written ``with many tears'' (2~Cor 2:1--4), and, on the most
common judgments, Galatians and perhaps the captivity notes that some
scholars assign to an unrecorded Ephesian imprisonment---Paul does say ``I
die daily'' and that he ``fought with beasts at Ephesus'' (1~Cor 15:31--32,
almost certainly figurative), and 2~Corinthians 1:8--10 recalls a mortal
danger in Asia that Acts does not narrate. Acts instead gives the
silversmiths' riot---``Great is Diana of the Ephesians'' (19:28)---as the
climax and departure-spring. Then Macedonia, three months in Greece
(Corinth, where Romans was written commending Phoebe of Cenchrae, Rom
16:1), and the deliberate last journey: back through Macedonia and Troas,
the farewell at Miletus to the Ephesian elders with its unique unwritten
saying of Jesus, ``It is a more blessed thing to give, rather than to
receive'' (Acts 20:35), and on toward Jerusalem for Pentecost with the
collection---against every warning: ``I am ready not only to be bound, but
to die also in Jerusalem, for the name of the Lord Jesus'' (21:13).

\begin{studybox}{The itineraries in brief (Acts' framework; DR spellings,
modern names where they differ)}
\small
\textbf{First journey} (Acts 13--14, c.~46--48): Antioch --- Seleucia ---
Cyprus (Salamis, Paphos) --- Perga --- Antioch of Pisidia --- Iconium ---
Lystra --- Derbe --- back through the same cities --- Attalia ---
Antioch.\par
\textbf{Second journey} (Acts 15:40--18:22, c.~49--52): Antioch --- Syria
and Cilicia --- Derbe, Lystra --- Phrygia and Galatia --- Troas ---
Neapolis --- Philippi --- Thessalonica --- Beroea --- Athens --- Corinth
(eighteen months) --- Cenchrae --- Ephesus (in passing) --- Caesarea ---
[Jerusalem] --- Antioch.\par
\textbf{Third journey} (Acts 18:23--21:16, c.~52--57): Antioch ---
Galatia and Phrygia --- Ephesus (about three years) --- Macedonia ---
Greece/Corinth (three months) --- Philippi --- Troas --- Assos ---
Mitylene --- Miletus --- Coos, Rhodes, Patara --- Tyre --- Ptolemais ---
Caesarea --- Jerusalem.\par
\textbf{Voyage in custody} (Acts 27--28, c.~59--60): Caesarea --- Sidon
--- the Lycian coast (Myra; DR ``Lystra'') --- Good-havens, Crete ---
storm --- Malta (three months) --- Syracusa (Syracuse) --- Rhegium ---
Puteoli --- Appii Forum and the Three Taverns --- Rome.
\end{studybox}
